\documentclass[../main.tex]{subfiles}
\graphicspath{{\subfix{../img/}}}
\begin{document}
	\subsection{Legge di Faraday-Lenz}
	Esperimenti condotti da faraday in inghilterra ed henry negli stati uniti nel 1821 misero in evidenza come magnetismo ed elettricità fossero strettamente connessi: un campo magnetico variabile nel tempo è ingrado di generare un campo elettrico.\\
	L’osservazione di correnti magneticamente indotte può essere effettuata prendendo una spira e un magnete (ed un amperometro) (figura \ref{fig:image26})
	\begin{figure}[h]
		
		\begin{subfigure}{0.33\textwidth}
			\includegraphics[scale=0.55]{185.png} 
			\caption{}
			\label{fig:im-185}
		\end{subfigure}
		\begin{subfigure}{0.33\textwidth}
			\includegraphics[scale=0.55]{186.png}
			\caption{}
			\label{fig:im-186}
		\end{subfigure}
		\begin{subfigure}{0.33\textwidth}
			\includegraphics[scale=0.55]{187.png}
			\caption{}
			\label{fig:im-187}
		\end{subfigure}
		
		\caption{}
		\label{fig:image26}
	\end{figure}
	Si registra una corrente nel circuito ogni volta che il magnete è posto in movimento nelle vicinanze della spira.\\
	Il verso della corrente cambia a seconda che il magnete venga \textbf{avvicinato od allontanato}.\\
	Il verso della corrente si inverte anche cambiando la polarità del magnete e gli effetti solla spira non cambiano se invece di muovere il magnete è la spira a muoversi (figura \ref{fig:image27})
	\begin{figure}[h]
		
		\begin{subfigure}{0.55\textwidth}
			\includegraphics[scale=0.55]{188.png} 
			\caption{}
			\label{fig:im-188}
		\end{subfigure}
		\begin{subfigure}{0.55\textwidth}
			\includegraphics[scale=0.55]{189.png}
			\caption{}
			\label{fig:im-189}
		\end{subfigure}
		
		\caption{}
		\label{fig:image27}
	\end{figure}
	Ogni volta che si ha un moto relativo tra spira e magnete si osserva la comparsa di una corrente $I$ \textbf{indotta} nel circuito.\\
	L’intensità della corrente che si osserva è proporzionale alla \textbf{rapidità della variazione}.\\\\
	Effetti analoghi si osservano qualora il magnete vengà sostituito da una sorgente di campo magnetico quale un conduttore percorso da corrente (\textbf{circuito primario}) (figura \ref{fig:image28})
	\begin{figure}[h]
		
		\begin{subfigure}{0.55\textwidth}
			\includegraphics[scale=0.55]{190.png} 
			\caption{}
			\label{fig:im-190}
		\end{subfigure}
		\begin{subfigure}{0.55\textwidth}
			\includegraphics[scale=0.55]{191.png}
			\caption{}
			\label{fig:im-191}
		\end{subfigure}
		
		\caption{}
		\label{fig:image28}
	\end{figure}
	Le correnti misurate si invertono invertendo la polarità della batteria.\\
	Ancora una volta ciò che “sembra contare” è il moto relativo di una spira rispetto all’altra.\\
	In nessuno dei casi descritti il \textbf{circuito (secondario)} è connesso ad alcuna batteria!!!
	$\Rightarrow$ La corrente indotta è \textbf{generata} da una forza elettromotrice $\mathscr{F}$ (f.e.m.) indotta.
	La f.e.m. $\mathscr{F}$ (volt) come nell’analogo della batteria è l’energia per unità di carica che deve essere fornita ai portatori di carica che costituiscono la corrente.\\\\
	Se con $R$ indichiamo la resistenza complessiva del circuito
	$$I = \frac{\mathscr{F}}{R}$$
	$\Rightarrow$ $I$ dipende da $R$ mentre la f.e.m. $\mathscr{F}$ è l’effetto principale del fenomeno che stiamo discutendo.\\
	Se ohmico.\\
	Infatti se interrompiamo il circuito in un punto ($R=\infty$) (figura \ref{fig:im-192})
	\begin{center}
		\begin{figure}[h]
			\centering
			\includegraphics[scale=0.55]{192.png}
			\caption{}
			\label{fig:im-192}
		\end{figure}
	\end{center}
	$$\boxed{I=0 \text{ ma } \mathscr{F} \neq 0}$$
	Si misura una f.e.m. $\mathscr{F}$ in tutto simile a quella che si ha ai capi di un generatore di tensione (la polarità dipende dal moto relativo!!)\\\\
	Per osservare una $\mathscr{F}$ (e quindi una $I$) può non essere necessaria una condizione di moto (figura \ref{fig:im-193})
	\begin{center}
		\begin{figure}[h]
			\centering
			\includegraphics[scale=0.55]{193.png}
			\caption{}
			\label{fig:im-193}
		\end{figure}
	\end{center}
	Si osserva una corrente indotta alla chiusura del circuito\\
	$\Rightarrow$ per tutto il tempo in cui la corrente nel primario si porta da $0$ ad un valore di regime.\\
	Si osserva nuovamente una corrente indotta all’apertura del circuito\\
	$\Rightarrow$ si osserva una corrente indotta negli intervalli di tempo in cui il \textbf{campo magnetico concatenato} con il circuito \textbf{varia nel tempo}.\\\\
	Queste osservazioni permisero a Faraday di formulare una legge generale per il fenomeno dell’\textbf{induzione elettromagnetica}.\\
	Faraday intuì che la f.e.m. indotta è proporzionale alla variazione nel tempo delle linee di forza che attraversano il circuito secondario\\
	$\Rightarrow$ la f.e.m. indotta è direttamente proporzionale alla derivata temporale del \textbf{flusso} del campo magnetico che attraversa il circuito.
	$$\boxed{\mathscr{F} = - \frac{d\phi_B}{dt}} \qquad \qquad \text{con } \phi_B = \int \vec{B} \cdot \vec{dS} \qquad \text{(Attraverso la superficie racchiusa nella spira)}$$
	È nota come la \textbf{legge di Faraday(-Neumann-Lenz) dell’induzione elettromagnetica}\\
	Nel S.I. $\pm \frac{Wb}{s} = \pm V$
	La legge di Faraday mostra che si osserva una f.e.m. indotta ogni qualvolta si ha una variazione del flusso di $\vec{B}$ concatenato con la spira.\\
	$\Rightarrow$
	\begin{itemize}
		\item Varia la superficie del circuito (deformazione del circuito nel tempo)
		\begin{itemize}
			\item Il campo $\vec{B}$ può essere sia uniforme che non uniforme
		\end{itemize}
		\item Il campo $\vec{B}$ (concatenato con la spira) varia nel tempo
		\begin{itemize}
			\item circuito fisso e sorgente in moto (ad eccezione del caso di un moto \textbf{traslatorio} in un campo uniforme)
			\item circuito in moto e sorgente fissa (ad eccezione del caso di un moto \textbf{traslatorio} in un campo uniforme)
			\item assenza di moto relativo ma con variazioni della corrente nel tempo che genera $\vec{B}$ o per presenza di variazioni locali della \textbf{permeabilità magnetica}
		\end{itemize}
		\item varia l’angolo $\boxed{\varphi}$ fra $\vec{B}$ e la normale $\vec{dA}$ alla superficie del circuito ($\Rightarrow$ es. rotazione)
	\end{itemize}
	Nella legge di Faradav è contenuto anche un segno negativo:
	$$\mathscr{F} = {\color{red}-} \frac{d\phi_B}{dt}$$
	Il suo significato è espresso da una legge nota come \textbf{Legge Di Lenz}, secondo la quale la corrente indotta in una spira ($\Rightarrow$ la polarità della f.e.m. indotta) chiusa conduttrice è sempre tale da opporsi alla variazione di flusso magnetico che la ha generata (esempi in figura \ref{fig:im-194}).
	\begin{center}
		\begin{figure}[h]
			\centering
			\includegraphics[scale=0.45]{194.png}
			\caption{}
			\label{fig:im-194}
		\end{figure}
	\end{center}
	La corrente indotta circolerà in modo da generare un campo $\vec{B}$ il cui contributo al flusso concatenato sarà \textbf{in opposizione} alla variazione che lo ha generato.\\\\
	La legge di Lenz traduce il \textbf{principio di conservazione dell’energia} : se la corrente indotta concorresse ad aumentare un flusso già di per se crescente si andrebbe nel senso di un aumento illimitato di corrente e quindi un generatore e di tensione in grado di erogare una potenza elettrica crescente non fornita dall’esterno.\\\\
	Attraverso un esempio (figura \ref{fig:im-195}) esaminiamo una origine fisica della legge di Faraday:
	Abbiamo il campo $\vec{B}$ uniforme
	\begin{center}
		\begin{figure}[h]
			\centering
			\includegraphics[scale=0.45]{195.png}
			\caption{La sbarretta slitta a sinistra e destra, conduce sui binari}
			\label{fig:im-195}
		\end{figure}
	\end{center}
	Circuito rettangolare metallico chiuso a sinistra da una sbarretta mobile con velocità $\vec{v} = \text{Costante}$\\\\
	Per la legge di Faraday:
	$$\phi_B = \int \vec{B} \cdot d\vec{s} = B \int ds = BLx \Rightarrow \frac{d\phi_B}{dt} = \frac{d(BLx)}{dt} = BL \frac{dx}{dt} = BLv$$
	$$\Rightarrow \mathscr{F} = BLv$$
	\begin{center}
		\begin{figure}[h]
			\centering
			\includegraphics[scale=0.45]{196.png}
			\caption{}
			\label{fig:im-196}
		\end{figure}
	\end{center}
	\begin{figure}[h]
		
		\begin{subfigure}{0.55\textwidth}
			\includegraphics[scale=0.55]{197.png} 
			\caption{}
			\label{fig:im-197}
		\end{subfigure}
		\begin{subfigure}{0.55\textwidth}
			\includegraphics[scale=0.55]{198.png}
			\caption{Nel Moto}
			\label{fig:im-198}
		\end{subfigure}
		
		\caption{}
		\label{fig:image29}
	\end{figure}
	Guardando la figura \ref{fig:im-197} notiamo che su ciascun elettrone della sbarretta agirà una forza:
	$$\vec{F}_B = -e\vec{v} \times \vec{B}$$
	$\Rightarrow$ in figura \ref{fig:im-198} si avrà un accumulo di cariche alle estremità della sbarretta.\\
	Il processo proseguirà fino all’equilibrio (figura \ref{fig:im-199})
	\begin{center}
		\begin{figure}[h]
			\centering
			\includegraphics[scale=0.45]{199.png}
			\caption{}
			\label{fig:im-199}
		\end{figure}
	\end{center}
	$$\Rightarrow F_E = F_B \Rightarrow -eE = -evB \Rightarrow E = vB$$
	$$\Rightarrow V_B - V_A = -\int_A^B \vec{E} \cdot d\vec{l} = \int_A^B E \, dl = E \int_A^B dl = EL$$
	$$\Rightarrow V_B - V_A = EL = vBL \equiv \mathscr{F}$$
	Il fenomeno della induzione elettromagnetica è quindi riconducibile in questo esempio alla forza di Lorenz.\\
	Il risultato è generalizzabile a tutti icasi in cui un circuito è in moto in un campo magnetico.\\\\
	Una ulterore considerazione sull’esempio appena trattato:\\
	Se nella sbarretta circola una corrente: 
	$$I = \frac{\mathscr{F}}{R} \Rightarrow I = \frac{BLv}{R}$$
	allora sulla sbarretta agirà una forza:
	$$\vec{F}_B = I \vec{L} \times \vec{B} \Rightarrow F_B = \frac{BLv}{R} L B = \frac{B^2 L^2 v}{R}$$
	La variazione del campo magnetico, da un lato genera la corrente $I$, dall’altro origina una \textbf{forza resistente} rispetto al moto della sbarretta (figura \ref{fig:im-200}).\\
	Si parla di \textbf{Attrito Elettromagnetico}.
	\begin{center}
		\begin{figure}[h]
			\centering
			\includegraphics[scale=0.45]{200.png}
			\caption{}
			\label{fig:im-200}
		\end{figure}
	\end{center}
	$$\Rightarrow$$ per mantenere la sbarretta in moto con una velocità $\vec{v}$ costante occorre applicare una forza esterna: 
	$$\vec{F} = -\vec{F}_B$$
	Spendendo una potenza: 
	$$P = \vec{F} \cdot \vec{v} = \frac{B^2 L^2 v^{\color{red}2}}{R}$$
	Ma essendo: 
	$$BLv = RI \Rightarrow P = \frac{\mathscr{F}^2}{R} = RI^2 = \underbrace{\mathscr{F} I}_{\text{potenza erogata da un generatore elettrico}}$$
	L’energia \textbf{meccanica} che si spende per mantenere in moto la sbarretta è \textbf{integralmente trasferita} al circuito come energia \textbf{elettrica}\\\\
	Il significato energetico della legge di Lenz è soddisfatto:\\
	Potenza Meccanica $\rightarrow$ Potenza Elettrica\\
	\textbf{Abbiamo realizzato un generatore di corrente}.\\\\
	Il processo inverso:
	Potenza Elettrica $\rightarrow$ Potenza Meccanica\\
	\textbf{È alla base del funzionamento dei motori elettrici}.
	\subsubsection{Spira in Moto in Un Campo Magnetico Uniforme}
	\begin{center}
		\begin{figure}[h]
			\centering
			\includegraphics[scale=0.45]{201.png}
			\caption{}
			\label{fig:im-201}
		\end{figure}
	\end{center}
	Facendo riferimento alla figura \ref{fig:im-201}:
	$$\phi_B = \int \vec{B} \cdot d\vec{s} = B \cos \vartheta \int ds = B A \cos \vartheta$$
	Ma se un agente esterno sta facendo ruotare la spira intorno ad un asse con una velocità angolare $\omega$ costante:
	$$\Rightarrow \vartheta = \vartheta(t) = \omega t$$
	$$\Rightarrow \phi_B = B A \cos \omega t$$
	Il flusso $\phi_B$ varierà nel tempo $\Rightarrow$ anche la f.e.m indotta (figura \ref{fig:im-202}):
	$$\mathscr{F} = - \frac{d\phi_B}{dt} = B A \omega \sin \omega t$$
	\begin{center}
		\begin{figure}[h]
			\centering
			\includegraphics[scale=0.45]{202.png}
			\caption{Tratteggiato è l’andamento di $\phi_B$}
			\label{fig:im-202}
		\end{figure}
	\end{center}
	Se nella spira circola una corrente $I$ $\Rightarrow$ agirà sulla spira un momento torcente: 
	$$\vec{\tau} = I \vec{A} \times \vec{B}$$
	\begin{center}
		\begin{figure}[h]
			\centering
			\includegraphics[scale=0.45]{203.png}
			\caption{}
			\label{fig:im-203}
		\end{figure}
	\end{center}
	Osservando quindi figura \ref{fig:im-203} notiamo che il momento tenderebbe a produrre una rotazione nel verso opposto alla rotazione della spira\\
	$\Rightarrow$ occorre compiere un lavoro meccanico per mantenere in moto la spira.\\
	\textbf{Abbiamo realizzato un Generatore Di Corrente Alternata}
	\subsubsection{Analogia con Legge di Ampere-Maxwell}
	La forza di Lorentz è responsabile della legge di Faraday \textbf{se il conduttore è in moto}\\
	Se il conduttore e la sorgente \textbf{sono fermi} non possiamo ricorrere alla forza di Lorentz.\\
	Eppure la presenza di una corrente in una spira al variare del flusso concatenato implica la presenza di una forza agente e quindi di un campo elettrico $\vec{E}$ nel conduttore.
	Essendo la f.e.m. $\mathscr{F}$ il lavoro per unità di carica compiuto dalla forza elettrica agente sui portatori di carica:
	$$\mathscr{F} = \frac{W}{q} = \oint \frac{\vec{F} \cdot d\vec{l}}{q} = \oint \vec{E} \cdot d\vec{l}$$
	$$\mathscr{F} = \oint \vec{E} \cdot d\vec{l}$$
	La f.e.m. indotta è legata all’insorgere nel conduttore di un campo elettrico $\vec{E}$ indotto\\
	$\Rightarrow$ è un campo elettrico non conservativo $\Rightarrow$ non è il risultato di una distribuzione statica di cariche per le quali $\oint \vec{E} \cdot d\vec{l} = 0$\\\\
	Se al posto della spira conduttrice si sostitviste un percorso chiuso:
	$$\boxed{\oint \vec{E} \cdot d\vec{l} = - \frac{d\phi_B}{dt}}$$
	$\Rightarrow$ La variazione del flusso di campo magnetico è \textbf{essa stessa} sorgente di un campo elettrico nello spazio.\\\\
	Divieve evidente l’analogia con la legge di Ampère-Maxwell (in assenza di correnti di conduzione concatenate):
	$$\oint \vec{E} \cdot d\vec{l} = \mu_0 \varepsilon_0 \frac{d\phi_E}{dt} \qquad \qquad \oint \vec{E} \cdot d\vec{l} = - \frac{d\phi_B}{dt}$$
	\textbf{Campi elettrici e magnetici variabili nel tempo si generano reciprocamente}.\\
	La sintesi sarà completata nella nozione di \textbf{campo elettromagnetico}.
	\subsection{Autoinduzione}
	La legge di Faraday $\mathscr{F} = -\frac{d\phi_B}{dt}$ è stata fino ad ora utilizzata per valutare la f.e.m. indotta a seguito di variazioni nel tempo del flusso $\phi_B$ concatenato con una spira conduttrice: \textbf{il campo magnetico era generato da una sorgente esterna}
	$\Rightarrow$ anche la spira stessa percorsa da corrente è essa stessa sorgente di un campo magnetico concatenato con la spira (figura \ref{fig:im-204})
	\begin{center}
		\begin{figure}[h]
			\centering
			\includegraphics[scale=0.45]{204.png}
			\caption{}
			\label{fig:im-204}
		\end{figure}
	\end{center}
	$\Rightarrow$ Il campo magnetico generato dalla spira stessa produce un flusso anche attraverso la spira 
	$$\Rightarrow \phi_B = \int \vec{B} \cdot d\vec{s} = \iint \frac{\mu_0 I}{4\pi} \overbrace{\frac{\vec{dl} \times \hat{r}}{r^2}}^{\text{Biot-Savart}} \cdot d\vec{s} = LI$$
	$$\Rightarrow \phi_B = LI$$
	L’\textbf{autoflusso} dipende dalla corrente $I$ e dalle caratteristiche geometriche del circuito.\\
	$L$ è detto \textbf{coefficiente di auto induzione} o \textbf{induttanza}.\\
	$\Rightarrow$ Per la legge di Faraday 
	$$\mathscr{F} = -\frac{d\phi_B}{dt} = -L\frac{dI}{dt}$$
	\begin{center}
		\begin{figure}[h]
			\centering
			\includegraphics[scale=0.45]{205.png}
			\caption{}
			\label{fig:im-205}
		\end{figure}
	\end{center}
	Facendo riferimento alla figura \ref{fig:im-205}: La f.e.m. indotta sarà in accordo con la legge di Lenz.\\
	Si parla di \textbf{auto induzione}: la variazione del flusso concatenato è originato dal circuito stesso.\\\\
	$\mathscr{F}$ È detta \textbf{f.e.m. auto indotta} o \textbf{forza contro elettro motrice}.\\
	Nel S.I. l’induitanza si misura in \textbf{Henry (H)}
	$$\pm H = \pm \frac{V \cdot S}{A}$$
	$$\Rightarrow \mu_0 = 4 \pi \cdot 10^{-7} \frac{T \cdot m}{A} = 4 \pi \cdot 10^{-7} {\color{red} \frac{H}{m}}$$
	\subsubsection{Induttanza $L$ Di Un Solenoide}
	All’interno del solenoide il campo magnetico è uniforme e parallelo all’asse del solenoide e pari a $B = \mu_0 n I$ in modulo.\\
	$\Rightarrow$ Ogni avvolgimento è approssimabile ad una spira piana di area $A$ attraverso la quale è concatentato il flusso $BA = \mu_0 \; n \: I \; A$\\
	$\Rightarrow$ Se $l$ è la lunghezza del solenoide il numero degli avvolgimenti e $nl$\\\\
	Il flusso totale concatenato è $\phi_B = nl \; \mu_0 \; n \; I \; A$\\
	$\Rightarrow \phi_B = \mu_0 \; n^2 \; l \; A \; I$ $\Rightarrow$ l'induttanza $L = \mu_0 n^2 l A = \mu_0 n^2 \underbrace{V}_{\text{Volume Solenoide}}$
	\subsection{Circuiti LR}
	Consideriamo un circuito che contenga un elemento induttivo (es un solenoide) [è detto induttore e sarà rappresentato da una induttanza $L$] in serie con una resistenza $R$.\\
	Per la legge di Faraday ad ogni variazione di corrente circolante nel circuito l’induttore si oppone producendo una f.e.m. autoindotta.\\\\
	In presenza di una corrente continua un induttore è "solo" un filo (se ideale $R=0$).
	\begin{center}
		\begin{figure}[h]
			\centering
			\includegraphics[scale=0.45]{206.png}
			\caption{}
			\label{fig:im-206}
		\end{figure}
	\end{center}
	Facendo riferimento alla figura \ref{fig:im-206} per $t=0$ il curcuito viene chiuso:\\
	La corrente del solenoide da $I=0$ inizierà ad aumentare $\Rightarrow I(t)$\\
	$\Rightarrow$ Nell'induttanza si genererà una f.e.m. autoindotta che si oppone all'aumento della corrente $I$.\\
	$\Rightarrow$ La f.e.m. indotta impedisce alla $I$ di aumentare bruscamente.\\
	$\Rightarrow$ La corrente nel circuito è quindi controllata da due f.e.m.:\\
	una costante $\mathscr{F}_0$ della batteria, ed una variabile $\mathscr{F}_L = -L \frac{dI}{dt}$, di segno opposto.\\
	$\Rightarrow$ Finchè $\mathscr{F}_L$ non si annulla, la corrente nel circuito sarà minore di $\frac{\mathscr{F}_0}{R}$.\\\\
	Applicando la legge delle Maglie: 
	$$(V_B - V_A) + (V_C - V_B) + (V_A - V_C) = 0$$
	$$\mathscr{F}_0 - IR - L \frac{dI}{dt} = 0$$
	Risolvendo:
	$$L \frac{dI}{dt} + RI = \mathscr{F}_0 \Rightarrow L dI = -R \left( I - \frac{\mathscr{F}_0}{R} \right) dt$$
	$$\Rightarrow \frac{dI}{I - \frac{\mathscr{F}_0}{R}} = -\frac{R}{L} dt \underbrace{\Rightarrow}_{\text{integrando}} \int_0^I \frac{dI}{I - \frac{\mathscr{F}_0}{R}} = \int_0^t -\frac{R}{L} dt$$
	$$\Rightarrow \ln \left( I - \frac{\mathscr{F}_0}{R} \right) - \ln \left( -\frac{\mathscr{F}_0}{R} \right) = -\frac{R}{L} t \Rightarrow I = \frac{\mathscr{F}_0}{R} \left( 1 - e^{-\frac{t}{\tau_L}} \right)$$ 
	Dove: $\tau_L = \frac{L}{R}$
	\begin{center}
		\begin{figure}[h]
			\centering
			\includegraphics[scale=0.45]{207.png}
			\caption{}
			\label{fig:im-207}
		\end{figure}
	\end{center}
	Osservando il grafico \ref{fig:im-207}:
	$$I = \frac{\mathscr{F}_0}{R} (1 - e^{-\frac{t}{\tau_L}}) \qquad t=0 \qquad I=0 \quad \Rightarrow \mathscr{F}_L = -\mathscr{F}_0$$
	$$t \to \infty \quad I \to \frac{\mathscr{F}_0}{R} \quad \text{ Valore di regime}$$
	$$\tau_L = \frac{L}{R} \quad \text{ È la costante di tempo del circuito}$$
	La corrente non aumenta istantaneamente ma esponenzialmente fino al valore di regime per la presenza di una corrente indottA $I_L = I(t) - I_{\infty} = \frac{-\mathscr{F}_0 \cdot e^{-\frac{t}{\tau}}}{R} = \frac{-\mathscr{F}_L}{R}$
	Analogamente se escluoiamo la batteria (figura \ref{fig:im-208})
	\begin{center}
		\begin{figure}[h]
			\centering
			\includegraphics[scale=0.45]{208.png}
			\caption{}
			\label{fig:im-208}
		\end{figure}
	\end{center}
	La corrente dal valore di regime $I_0$ diminuirà\\
	$\Rightarrow$ Nell’induttanza si genererà una f.e.m. $\mathscr{F}_L$ autoindotta che si oppone alla diminuzione della corrente $I$\\
	$\Rightarrow$ Impedisce alla $I$ di diminuire bruscamente.\\
	$I$ continua a circolare sostenuta dA $\mathscr{F}_L$:
	$$-IR - L \frac{dI}{dt} = 0$$
	$$\Rightarrow \frac{dI}{I} = -\frac{R}{L} dt \Rightarrow \int_{I_0}^{I} \frac{dI}{I} = \int_{0}^{t} -\frac{R}{L} dt \Rightarrow \ln \frac{I}{I_0} = -\frac{R}{L} t$$
	\begin{center}
		\begin{figure}[h]
			\centering
			\includegraphics[scale=0.45]{209.png}
			\caption{}
			\label{fig:im-209}
		\end{figure}
	\end{center}
	Osservando il grafico \ref{fig:im-209}:
	$$t=0 \qquad I=I_0 \quad \Rightarrow \mathscr{F}_L = RI_0 \quad \text{Anche se la batteria è esclusa}$$
	$$t \to \infty \quad I \to 0$$
	Se esclusa la batteria circola una corrente $I$\\
	$\Rightarrow$ Nella resistenza i portatori di carica perderanno energia per effetto Joule dissipando una potenza $P = RI^2$\\\\
	Da dove proviene questa energia?
	\subsubsection{Densità di Energia Magnetica}
	\textbf{L'energia deve essere immagazzinata nell'induttore}.\\
	Nel circuito LR connesso alla batteria:
	$$\mathscr{F}_0 = RI + L \frac{dI}{dt}$$
	E multiplicando per $I$ ambo i membri:
	$$\underbrace{\mathscr{F}_0 I}_{\text{Potenza erogata dalla batteria}} = \underbrace{RI^2}_{\text{Potenza dissipata dall'effetto Joule della resistenza}} + LI \underbrace{\frac{dI}{dt}}_{\frac{dU_B}{dt}}$$
	Calcoliamo così il \textbf{lavoro per unità di tempo contro la f.e.m. indotta}.\\
	Per la conservazione dell’energia $U_B$ deve essere \textbf{l’energia immagazzinata nell’induttore}:
	$$\frac{dU_B}{dt} = LI \frac{dI}{dt} \Rightarrow \int_0^{U_B} dU_B = \int_0^I LI \, dI \Rightarrow \boxed{U_B = \frac{1}{2}LI^2}$$
	Una situazione analoga era stata osservata nel corso del processo di carica (e scarica) di un condensatore ($U_E = \frac{1}{2} \frac{Q^2}{C}$)\\\\
	Se l’induttore è un solenoide di area $A$ e lunghezza $l$:\\
	$L = \mu_0 n^2 Al$ $\quad$ e $\quad$ $B = \mu_0 n I \Rightarrow I = \frac{B}{\mu_0 n}$
	$$\Rightarrow U_B = \frac{1}{2}LI^2 = \mu_0 n^2 Al \left(\frac{B}{\mu_0 n}\right)^2 = {\color{red} \frac{1}{2} \frac{B^2}{\mu_0} Al}$$
	Se introduciamo la \textbf{densità di energia} $\mu_0 = \frac{U_0}{Al} \rightarrow \text{Volume Solenoide}$
	$$\Rightarrow \boxed{\mu_0 = \frac{1}{2} \frac{B^2}{\mu_0}} \qquad \textbf{L’energia È Immagazzinata Nel Campo Magnetico}$$
	Ha validità \textbf{generale}. \\
	È valida per ogni regione dello spazio in cui esiste un campo magnetico.\\
	Analogamente a quanto osservato per il campo elettrico dove $\mu_E = \frac{1}{2} \varepsilon_0 E^2$
	\subsection{Mutua Induzione e Trasformatori}
	Una variazione nell’intensità della corrente circolante in una bobina può far variare il flusso concatenato con un circuito presente nelle vicinanze (e viceversa) $\Rightarrow$ si parla di \textbf{Mutua Induzione}.
	\begin{center}
		\begin{figure}[h]
			\centering
			\includegraphics[scale=0.45]{210.png}
			\caption{}
			\label{fig:im-210}
		\end{figure}
	\end{center}
	Osserviamo figura \ref{fig:im-210}, se indichiamo con $\phi_{21}$ il flusso del campo magnetico prodotto dalla bobina $1$ e concatenato con la bobina $2$\\
	$\Rightarrow \phi_{21} \propto I_1$\\
	La costante di proporzionalità è detta \textbf{Coefficiente Di Mutua Induzione}\\
	$\Rightarrow \phi_{21} = M_{21} I_1$ e analogamente $\phi_{12} = M_{12} I_2$\\\\
	Per la legge di Faraday:\\
	$\mathscr{F}_{21} = - \frac{d\phi_{21}}{dt} = - M_{21} \frac{dI_1}{dt}$ e analogamente $\mathscr{F}_{12} = - \frac{d\phi_{12}}{dt} = - M_{12} \frac{dI_2}{dt}$\\
	Il coefficiente di mutua induzione dipende dalle caratteristiche geometriche delle bobine e dalla loro posizione relativa.\\\\
	Si può dirostrare che $M_{21} = M_{12} = M$\\
	$\Rightarrow \mathscr{F}_2 = - M \frac{dI_1}{dt}$ $\quad$ e $\quad$ $\mathscr{F}_1 = - M \frac{dI_2}{dt}$\\
	$\Rightarrow$ La f.e.m. indotta in un circuito e sempre proporzionale alla velocità di variazione della corrente nell’altro circuito.
	\subsection{Circuiti LC}
	\subsubsection{Circuito RC}
	\begin{center}
		\begin{figure}[h]
			\centering
			\includegraphics[scale=0.45]{211.png}
			\caption{}
			\label{fig:im-211}
		\end{figure}
	\end{center}
	Raffigurato in figura \ref{fig:im-211}:
	$$I = \underbrace{I_0}_{\frac{\Delta V}{R}} e^{-\frac{t}{\tau_R}} \qquad \tau_R = RC$$
	L’energia $U_B = \frac{1}{2} \frac{Q^2}{C}$ immagazzinata nel condensatore viene dissipata sottoforma di calore nella resistenza $R$.
	\subsubsection{Circuito RL}
	\begin{center}
		\begin{figure}[h]
			\centering
			\includegraphics[scale=0.45]{212.png}
			\caption{}
			\label{fig:im-212}
		\end{figure}
	\end{center}
	Raffigurato in figura \ref{fig:im-212}:
	$$I = I_0 \, e^{-\frac{t}{\tau_L}} \qquad \tau_L = \frac{L}{R}$$
	L’energia $U_B = \frac{1}{2} L I_0^2$ immagazzinata nell’induttore viene dissipata sottoforma di calore nella resistenza $R$.
	\subsubsection{Circuito LC}
	Consideriamo ora cosa accade in un circuito in cui siamo presenti una capacità $C$ ed una induttanza $L$ e nessuna resistenza. (figura \ref{fig:im-213})
	\begin{center}
		\begin{figure}[h]
			\centering
			\includegraphics[scale=0.45]{213.png}
			\caption{Circuito LC}
			\label{fig:im-213}
		\end{figure}
	\end{center}
	Alla chiusura del circuito il condensatore inizierà a scaricarsi attraverso l’induttanza.\\
	Nel circuito circolerà una corrente $I = \frac{dQ}{dt}$ \textbf{variabile nel tempo}.\\
	$\Rightarrow$ Avremo una corrente circolante nell’induttore \\
	$\Rightarrow$ Stiamo immagazzinando energia magnetica $U_B = \frac{1}{2} L I^2$ nell’induttore a spese dell’energia del condensatore $U_B = \frac{1}{2} \frac{Q^2}{C}$.
	\begin{center}
		\begin{figure}[h]
			\centering
			\includegraphics[scale=0.45]{214.png}
			\caption{Circuito LC Chiuso}
			\label{fig:im-214}
		\end{figure}
	\end{center}
	Guardando attentamente figura \ref{fig:im-214} osservaimo come:
	$$(V_A - V_B) + (V_B - V_A) = 0$$
	$$\frac{Q}{C} - L \frac{dI}{dt} = 0$$
	Il condensatore si sta scaricandO $\Rightarrow I = -\frac{dQ}{dt}$
	$$\Rightarrow L \frac{d^2Q}{dt^2} = -\frac{Q}{C} \quad \Rightarrow \frac{d^2Q}{dt^2} = -\frac{1}{Lc} Q$$
	L’equazione differenziale è la “stessa” per il sistema \textbf{"massa-molla"} (figura \ref{fig:im-215}).
	\begin{center}
		\begin{figure}[h]
			\centering
			\includegraphics[scale=0.45]{215.png}
			\caption{Sistema Massa-Molla}
			\label{fig:im-215}
		\end{figure}
	\end{center}
	$$M \frac{d^2x}{dt^2} = -Kx \Rightarrow \frac{d^2x}{dt^2} = -\frac{K}{M}x$$
	$$\begin{gathered} 
		x \leftrightarrow Q \\ 
		v_x \leftrightarrow I \\ 
		K \leftrightarrow \frac{1}{C} \\ 
		M \leftrightarrow L 
		\end{gathered}$$
	Avente come soluzione $x = x_H \cos(\omega t + \varphi)$ con $\omega = \sqrt{\frac{K}{M}}$
	$$\Rightarrow$$ La carica sul condensatore nei circuiti $$LC$$ varierà sinusoidalmente secondo la relazione:
	$$Q = Q_M \cos(\omega t + \varphi) \quad \text{ con } \quad \omega = \frac{1}{\sqrt{LC}}$$
	$Q_M$ è la carica massima sul condensatore, $\omega$ la pulsazione della \textbf{oscillazione}, $$\varphi$$ costante di fase.\\
	Se al tempo $t=0$ $\quad$ $Q=Q_M \Rightarrow \varphi=0$ E $Q=Q_M \cos \omega t$\\
	Anche la corrente avrà un andamento oscillante:
	$$I = -\frac{dQ}{dt} = \omega Q_M \sin \omega t = I_M \sin \omega t$$
	Abbiamo che in un circuito LC ideale sia la corrente che la carica varieranno sinusoidamente nel tempo con pulsazione $\frac{1}{\sqrt{LC}}$
	$$\Rightarrow y = \frac{1}{2\pi} \frac{1}{\sqrt{LC}}$$
	$$Q(t) \quad \text{e} \quad I(t)$$ sono sfasate di $90^\degree$ (figura \ref{fig:image30})
	\begin{figure}[h]
		
		\begin{subfigure}{0.55\textwidth}
			\includegraphics[scale=0.55]{216.png} 
			\caption{}
			\label{fig:im-216}
		\end{subfigure}
		\begin{subfigure}{0.55\textwidth}
			\includegraphics[scale=0.55]{217.png}
			\caption{}
			\label{fig:im-217}
		\end{subfigure}
		
		\caption{}
		\label{fig:image30}
	\end{figure}
	Quando il condensatore è scarico, la corrente che circola nel circuito è massima\\
	$\Rightarrow$ Sarà massima l’energia immagazzinata nell’induttanza L in assenza di resistenze \textbf{non ci sarà dissipazione di energia} sotto forma di calore.\\
	$\Rightarrow$ L’energia sarà immagazzinata unicamente sottoforma di energia elettrica:
	$$\text{Nel condensatore carico: } U_E = \frac{1}{2} \frac{Q^2}{C} = \frac{1}{2} \frac{Q_M^2}{C} \cos^2 \omega t$$
	E di energia magnetica:
	$$\text{Nell'induttanza percorsa da corrente: } U_B = \frac{1}{2} L I^2 = \frac{1}{2} L I_M^2 \sin^2 \omega t$$
	L’energia totale è data da $\boxed{U = U_B + U_E}$ \textbf{Costante}\\\\
	Essendo $I_M = \omega Q_M \Rightarrow I_M = \frac{1}{\sqrt{LC}} Q_M \Rightarrow I_M^2 = \frac{Q_M^2}{LC} \Rightarrow \frac{1}{2} L I_M^2 = \frac{1}{2} \frac{Q_M^2}{C}$\\
	$\Rightarrow U = U_E + U_B = \underbrace{\frac{1}{2} L I_M^2}_{\equiv \rightarrow} \sin^2 \omega t \cdot \underbrace{\frac{1}{2} \frac{Q_M^2}{C}}_{\leftarrow\equiv} \cos^2 \omega t = \frac{1}{2} \frac{Q_M^2}{C} = \frac{1}{2} L I_M^2$
	$\Rightarrow$ L’energia elettromagnetica totale nel circuito rimane costante trasformandosi continuamente da energia elettrica nel condensatore ad energia magnetica nell’industanza e viceversa (figura \ref{fig:im-218})
	\begin{center}
		\begin{figure}[h]
			\centering
			\includegraphics[scale=0.45]{218.png}
			\caption{}
			\label{fig:im-218}
		\end{figure}
	\end{center}
	Segno di (figura \ref{fig:im-219})
	\begin{center}
		\begin{figure}[h]
			\centering
			\includegraphics[scale=0.45]{219.png}
			\caption{}
			\label{fig:im-219}
		\end{figure}
	\end{center}
	Convenzione in figura \ref{fig:im-219}: $Q$ sulla faccia superiore del condensatore $I > 0$ se in verso antiorario.
\end{document}